\chapter*{Preface}
\addcontentsline{toc}{chapter}{Preface}
\textit{For the best interests of my daughter, Rea.}

\newpage

Parenting is a profound journey a human being can undertake. 
It asks us to show up fully---not just as caregivers and providers, but as mirrors, mentors, and companions to the child who looks to us to make sense of the world.

Yet this path can feel overwhelming. 
In the heat of a tantrum, in the exhaustion of  a sleepless night, or in the middle of an argument, we may find ourselves reacting in ways that don't reflect the parent we want to be.

This book is written for those moments.
\textit{The Empathetic Parent} is not trying to make a child `obedient' or getting perfect `A's or fulfilling unfulfilled dreams of a parent. 
It is a practical guide grounded in decades of research in developmental psychology, neuroscience, and family systems therapy. 

\textbf{Empathy is all you need.} 
The core thesis of this book is that: a child grows into a healthy, wise, and resilient adult with parenting that emphasizes emotion awareness, connection, and love, compared to method that imposes emotion dismissal, punishment, and control.

The pages that follow are organized as a journey: from understanding what empathy is and why it matters, to the everyday skills of active listening and emotion coaching, to navigating the particular challenges of each developmental stage.
Each chapter starts with a quote from a domain expert.
The main body details the information supported by scientific evidence and research with citations and addresses keywords with bold font.
`Key Insight' by end of each section provides a TLDR that summarizes the section.
One may use this to skip to the sections of interest.
The `Practice Exercise' section includes multiple choice questions to detect gaps in understanding.
`Reflection Questions' section poses open-ended questions to encourage reflection.
 
No matter where you are in your parenting journey---whether you are expecting your first child, raising a curious toddler, navigating the turbulent waters of adolescence, or hoping to repair or deepen a relationship with a grown child---this book is here with you.
